% Consolidated 17-electro-conductors from 309dbbd2c1c3ab17ea53128b49313f224ecea5a5; original TeX fragments preserved.

% === BEGIN 17-electro-conductors-054-065.tex blob 62577b852cd7edeb9818954f7bd1f96997ed8113 ===
% Irodov Part III §3.2, original 54--65, PDF pp.108--109,
% printed pp.99--100. Image-charge signs checked; conductor assumed ideal.
\electroconductproblem{54}{99}
一个小球由劲度系数为 $k_s$ 的绝缘弹性细线
悬挂在无限大水平导体平面上方。
小球充电后比原平衡位置向下降
$x$，此时距导体平面 $\ell$。
求小球电荷模。
\exsolution
无限导体平面用镜像法，
在平面下方距平面同为 $\ell$ 的位置
放置大小相同、异号的镜像电荷。
两电荷相距 $2\ell$，
作用于真小球的镜像吸引力为
\[
F_e=\frac{q^2}{4\pi\varepsilon_0(2\ell)^2}
=\frac{q^2}{16\pi\varepsilon_0\ell^2}.
\]
充电前后球的重力相同；
新增加的弹性线拉力为 $k_sx$，
故 $F_e=k_sx$。
因此
\[
\boxed{|q|=4\ell\sqrt{\pi\varepsilon_0k_sx}}.
\]
\excheck $x$ 若以 cm 给出，
代数值时须先换算成 m；
结果与平面是否带净电荷无关的前提是
理想接地或足够大的导体平面。

\electroconductproblem{55}{99}
点电荷 $q$ 位于距无限大接地导体平面
$\ell$ 处。
求把该电荷缓慢移到无穷远需要做的外功。
\exsolution
镜像电荷为 $-q$，位于平面另一侧
对称位置。真电荷距平面 $z$ 时，
受到指向平面的吸引力
$F_z=-kq^2/(2z)^2=-q^2/(16\pi\varepsilon_0z^2)$。
为准静态地从 $z=\ell$ 移至无穷远，
外力必须做
\[
W_{\rm ext}=-\int_\ell^\infty F_z\,\mathrm dz
=\frac{q^2}{16\pi\varepsilon_0}
\int_\ell^\infty z^{-2}\,\mathrm dz.
\]
故
\[
\boxed{W_{\rm ext}=\frac{q^2}{16\pi\varepsilon_0\ell}}.
\]
\excheck 不应把真电荷与镜像电荷距离当作固定；
且镜像电荷随真电荷移动，
直接用“两真实点电荷的势能”会有
多出一倍的计数问题。

\electroconductproblem{56}{99}
两个真点电荷 $+q,-q$ 在同一高度
$h=\ell/2$、相互水平距离为 $\ell$，
二者均位于无限接地导体平面上方。
求：(a) 每个真电荷受力模；
(b) 两电荷中点处电场强度模。
\exsolution
取平面为 $z=0$，真电荷位置为
$(0,h),(\ell,h)$。
对应镜像分别为 $-q$ 于 $(0,-h)$、
$+q$ 于 $(\ell,-h)$。
(a) 真正的 $+q$ 受三项作用：
真 $-q$ 的水平吸引
$+kq^2/\ell^2$，
自身镜像的向下吸引
$-kq^2/\ell^2$，
另一镜像 $+q$ 的斜向排斥
$kq^2/(2\ell^2)$、方向 $(-1,1)/\sqrt2$。
合力分量互为相反数：
\[
F_x=\frac{kq^2}{\ell^2}
\left(1-\frac{1}{2\sqrt2}\right),\quad F_z=-F_x.
\]
所以
\[
\boxed{F=\frac{(2\sqrt2-1)q^2}
{8\pi\varepsilon_0\ell^2}}.
\]
(b) 中点两真电荷电场沿连线同向叠加，
贡献 $8kq/\ell^2$；
两个镜像对该方向的合贡献为
$-8kq/(5\sqrt5\,\ell^2)$，
竖直分量相消。故
\[
\boxed{E=
\frac{2|q|}{\pi\varepsilon_0\ell^2}
\left(1-\frac1{5\sqrt5}\right)}.
\]
\excheck 所有镜像仅为解边界值的数学工具，
不是真实可自由移动的第三、第四个点电荷。

\electroconductproblem{57}{99}
两个无限大、相互垂直的接地导体半平面形成直角，
一点电荷 $q$ 位于两平面之间，
距每个平面均为 $\ell$。
求该电荷所受合电场力模。
\exsolution
在直角区域坐标中取真电荷 $(\ell,\ell)$。
两个单次反射镜像电荷为 $-q$，
位置分别在 $(-\ell,\ell)$、
$(\ell,-\ell)$；
双次反射镜像为 $+q$，
位置在 $(-\ell,-\ell)$。
前两个分别沿 $-x,-y$ 方向吸引，
每个力模为 $kq^2/(4\ell^2)$。
第三个沿 $(+x,+y)$ 排斥，
力模为 $kq^2/(8\ell^2)$，
每分量为该值除以 $\sqrt2$。
所以
\[
F_x=F_y=
-\frac{kq^2}{4\ell^2}
+\frac{kq^2}{8\sqrt2\,\ell^2}.
\]
合力模
\[
\boxed{F=
\sqrt2|F_x|=
\frac{(2\sqrt2-1)q^2}
{32\pi\varepsilon_0\ell^2}}.
\]
\excheck 必须有\emph{三个}镜像，
且双次反射镜像与真电荷同号；
漏掉它会造成错误的系数。

\electroconductproblem{58}{99}
点电偶极矩为 $\bm p$ 的偶极子
距接地无限大导体平面 $\ell$，
偶极矩垂直于导体表面。
求作用在偶极子上的力模。
\exsolution
垂直导体平面的偶极子的镜像偶极矩
与原偶极矩\emph{同向}，
两偶极子中心距为 $2\ell$。
在两同轴同向偶极子的轴线上，
相互作用势能函数的距离部分为
$U=-2kp^2/d^3$。
对\emph{真偶极子在镜像场中的力}，
在固定镜像构造下求其所受局部场梯度：
\[
F=\frac{6kp^2}{d^4}\bigg|_{d=2\ell}
=\frac{6kp^2}{16\ell^4}.
\]
因此
\[
\boxed{F=\frac{3p^2}{32\pi\varepsilon_0\ell^4}}.
\]
方向朝向导体平面。
\excheck 能量法处理镜像随距离改变时，
镜像系统能量要乘 $1/2$；
直接用 $\bm F=(\bm p\cdot\nabla)\bm E_{\rm image}$
不会混淆这一因子。

\electroconductproblem{59}{99}
点电荷 $q$ 位于无限接地导体平面上方
高度 $\ell$。
求导体表面，距电荷垂足水平距离为 $r$
处的感应面电荷密度，并求总感应电荷。
\exsolution
取导体为 $z=0$，镜像 $-q$ 在 $z=-\ell$。
在平面上方，
两者的法向电场分量叠加为
\[
E_z(0^+,r)=
-\frac{q\ell}{2\pi\varepsilon_0
(r^2+\ell^2)^{3/2}}.
\]
导体内部 $\bm E=0$，
表面法向边界条件
$\sigma=\varepsilon_0E_z(0^+,r)$，
于是
\[
\boxed{\sigma(r)=
-\frac{q\ell}{2\pi(r^2+\ell^2)^{3/2}}}.
\]
积分整个平面：
\[
Q_{\rm ind}
=\int_0^\infty\sigma(r)2\pi r\,\mathrm dr
=-q\ell\int_0^\infty
\frac{r\,\mathrm dr}{(r^2+\ell^2)^{3/2}}
=\boxed{-q}.
\]
\excheck 对正点电荷，感应面电荷密度应全为负；
式子积分总量为 $-q$ 是重要自校。

\electroconductproblem{60}{99}
一根无限长线电荷密度 $\lambda$ 的细线，
平行于无限接地导体平面，距平面 $\ell$。
求 (a) 每单位长度所受电场力模；
(b) 平面上距细线所在垂直投影平面
横向距离 $x$ 的感应面电荷密度。
\exsolution
用镜像细线，线密度为 $-\lambda$，
平行放在平面另一侧，距真线 $2\ell$。
两无限直线之间每单位长度的静电作用力
\[
\boxed{\frac FL=
\frac{\lambda^2}{2\pi\varepsilon_0(2\ell)}
=\frac{\lambda^2}{4\pi\varepsilon_0\ell}}.
\]
方向朝向导体平面。
(b) 设导体面为 $z=0$，
真线位于 $(x,z)=(0,\ell)$，
镜像线位于 $(0,-\ell)$。
面上法向场
\[
E_z(0^+,x)
=-\frac{\lambda\ell}{\pi\varepsilon_0(x^2+\ell^2)}.
\]
故
\[
\boxed{\sigma(x)=
-\frac{\lambda\ell}{\pi(x^2+\ell^2)}}.
\]
\excheck 对正 $\lambda$，感应面电荷为负，
且沿每单位线长积分
$\int_{-\infty}^\infty\sigma(x)\,\mathrm dx=-\lambda$。

\electroconductproblem{61}{100}
半无限长直细线垂直于无限接地导体平面，
其近端距平面为 $\ell$，
线电荷密度为 $\lambda$。
设该线在平面上的垂足为 O，
求 (a) O 点感应电荷面密度；
(b) 与 O 距离 $r$ 处的分布。
\exsolution
设真线沿 $z\ge\ell$ 延伸，
每个高度 $z$ 的电荷元
$\mathrm dq=\lambda\,\mathrm dz$，
其镜像在 $-z$。
由第59题的点电荷镜像诱导密度叠加：
\[
\sigma(r)=
-\frac{\lambda}{2\pi}
\int_\ell^\infty
\frac{z\,\mathrm dz}{(r^2+z^2)^{3/2}}.
\]
积分原函数为
$-1/\sqrt{r^2+z^2}$，
因此
\[
\boxed{\sigma(r)=
-\frac{\lambda}{2\pi\sqrt{r^2+\ell^2}}}.
\]
在 $r=0$：
\[
\boxed{\sigma(0)=-\frac{\lambda}{2\pi\ell}}.
\]
\excheck 本题的半无限线总电荷无限，
因而不要求整个平面的感应电荷积分有限；
面密度随 $r$ 缓慢衰减。

\electroconductproblem{62}{100}
半径 $R$ 的均匀带电细圆环总电荷 $q$，
其所在平面平行于无限接地导体平面，
且距离为 $\ell$。
求 (a) 轴线投影点处的感应面电荷密度；
(b) 环心处的电场强度模和电势。
\exsolution
真实环在 $z=\ell$，
镜像为带 $-q$ 的相同圆环在 $z=-\ell$。
(a) 对底部投影点，
两圆环在导体平面上的法向场叠加，
\[
E_z(0)=
-\frac{q\ell}{2\pi\varepsilon_0(R^2+\ell^2)^{3/2}},
\]
所以
\[
\boxed{\sigma(0)=
-\frac{q\ell}{2\pi(R^2+\ell^2)^{3/2}}}.
\]
(b) 环心处真实圆环的自场为0，
镜像环位于轴向距离 $2\ell$ 处，
其场模为
\[
\boxed{E=
\frac{|q|(2\ell)}
{4\pi\varepsilon_0(R^2+4\ell^2)^{3/2}}}.
\]
电势由真实环及镜像环相加：
\[
\boxed{\varphi_{\rm center}=
\frac q{4\pi\varepsilon_0}
\left(\frac1R-\frac1{\sqrt{R^2+4\ell^2}}\right)}.
\]
\excheck 电场方向指向导体面（当 $q>0$）；
电势是标量总和，必须计入真实环自电势。

\electroconductproblem{63}{100}
一颗总电荷为零、半径为 $R$ 的
\emph{孤立}导体球，
球心外距 $\ell>R$ 处有点电荷 $q$。
求导体球的平衡电势。
\exsolution
接地导体球的镜像解是：
在距球心 $R^2/\ell$ 处的镜像电荷
$q'=-qR/\ell$，
它与外部点电荷共同保证球面电势为0。
现在球要求总电荷为0而非接地，
需在球心再加补偿镜像
$q''=-q'=qR/\ell$。
由于球心点电荷使球面各处电势
增加相同常数
$kq''/R$，
外部点电荷与偏心镜像在球面仍恰好抵消，
故
\[
\boxed{\varphi_{\rm sphere}
=\frac{q''}{4\pi\varepsilon_0R}
=\frac{q}{4\pi\varepsilon_0\ell}}.
\]
\excheck 球电势不依赖半径 $R$，
但诱导面电荷\emph{分布}依赖 $R$；
本题的球是孤立无净电荷，
不能直接当作接地导体球。

\electroconductproblem{64}{100}
一颗半径为 $R_1,R_2$（$R_1<R_2$）
的中空导体球壳原本不带净电荷。
点电荷 $q$ 位于球腔内部，
距球心 O 为 $r<R_1$。
求球心 O 处电势。
\exsolution
真实点电荷在 O 处产生电势 $kq/r$。
为使金属内部场为零，
内球面出现总感应电荷 $-q$，
外球面出现总感应电荷 $+q$。
尽管内球面的电荷分布可能不均匀，
全部位于距 O 恰为 $R_1$ 的球面上，
故其在 O 处总电势为
$k(-q)/R_1$。
外球面均匀总电荷 $+q$ 在 O 的电势为
$kq/R_2$。
因此
\[
\boxed{\varphi(O)=
\frac q{4\pi\varepsilon_0}
\left(\frac1r-\frac1{R_1}+\frac1{R_2}\right)}.
\]
\excheck 内表面\emph{不是}均匀感应电荷，
但对球心电势，每个表面元距离相同，
故只需其总量。

\electroconductproblem{65}{100}
两个同心薄导体球，
内球半径为 $a$、带电 $q_1>0$，
外球半径为 $b>a$。
应给外球加多少总电荷 $q_2$，
才能让内球电势为零？
求各区间电势 $\varphi(r)$ 并描述图线。
\exsolution
内球表面电势为
$\varphi(a)=k(q_1/a+q_2/b)$，
故令其为0得到
\[
\boxed{q_2=-\frac ba q_1}.
\]
静电平衡时两球外电势按包围电荷求：
$r\le a$ 内球内部为常势零；
$a\le r\le b$，
由内球场积分并使用 $\varphi(a)=0$：
$\varphi(r)=kq_1(1/r-1/a)$；
$r\ge b$，
两球等效总电荷
$q_1+q_2=q_1(1-b/a)$，
所以
\[
\boxed{\varphi(r)=\frac{q_1}{4\pi\varepsilon_0}
\begin{cases}
0,&0\le r\le a,\\
\dfrac1r-\dfrac1a,&a<r<b,\\
\dfrac{1-b/a}{r},&r\ge b.
\end{cases}}
\]
电势图像在 $r=a$ 从0起向负值下降，
到 $r=b$ 后以负的 $1/r$ 形式趋向0，
且在两个金属表面电势连续。
\excheck 外球带电绝对值
$|q_2|>q_1$，故远处总电荷为负；
题目并未要求整个系统中性。
% === END 17-electro-conductors-054-065.tex ===

% === BEGIN 17-electro-conductors-066-077.tex blob 9e18cbe36655028b143ae4006589b68b5b7e1723 ===
% Irodov Part III §3.2, problems 66--77, PDF pp.109--110,
% printed pp.100--101. Fig.3.8 and Fig.3.9 checked against page images.
\electroconductproblem{66}{100}
图3.8中四块相邻大金属板
从上到下编号1、2、3、4，
相邻间距均为 $d$。
外侧板1、4用导线连通，
内侧2、3分别接电源正负端，
其电势差为 $\Delta\varphi>0$。
求 (a) 相邻空隙电场；
(b) 每块板的\emph{总面电荷密度}。
\exsolution
记 $\varphi_2-\varphi_3=\Delta\varphi$，
外侧共同电势
$\varphi_1=\varphi_4=\varphi_0$。
外侧板总电荷守恒且起初无净电荷，
故两外板净面密度之和为0：
\[
\frac{\varepsilon_0}{d}(\varphi_0-\varphi_2)
+\frac{\varepsilon_0}{d}(\varphi_0-\varphi_3)=0
\quad\Longrightarrow\quad
\varphi_0=\frac{\varphi_2+\varphi_3}{2}.
\]
以向下为正方向，
各间隙场分量为
$E_{i,i+1}=(\varphi_i-\varphi_{i+1})/d$：
\[
\boxed{E_{12}=-\frac{\Delta\varphi}{2d},\quad
E_{23}=\frac{\Delta\varphi}{d},\quad
E_{34}=-\frac{\Delta\varphi}{2d}}.
\]
各板总面电荷密度等于相邻表面电荷之和：
\[
\boxed{\sigma_1=-\frac{\varepsilon_0\Delta\varphi}{2d},
\quad \sigma_2=\frac{3\varepsilon_0\Delta\varphi}{2d},
\quad \sigma_3=-\frac{3\varepsilon_0\Delta\varphi}{2d},
\quad \sigma_4=\frac{\varepsilon_0\Delta\varphi}{2d}}.
\]
\excheck $\sigma_1+\sigma_4=0$、
$\sigma_2+\sigma_3=0$，
且中间间隙场强大小为两边的2倍。

\electroconductproblem{67}{100}
两块无限大平行导体板1、2相距 $\ell$。
一个点电荷 $q$ 处于板间，
距板1的距离为 $x$（$0<x<\ell$）。
求两板上\emph{总感应电荷}。
\exsolution
点电荷的场可用沿板平行平面的傅里叶分解、
镜像级数或 Green 函数处理。
计算\emph{总}感应电荷时，
不用求详细面密度：
设点电荷距板1为 $x$，
将点电荷想象沿平行板平面均匀铺展，
由于该等效化只改变表面电荷分布而不改变
两板\emph{总量}，
法向静电场在两段间距
$x$ 与 $\ell-x$ 上为分段常值。
导体板电势相同，
两段电场电势降必须相抵，
得到
\[
\boxed{Q_1=-q\,\frac{\ell-x}{\ell}},
\qquad
\boxed{Q_2=-q\,\frac{x}{\ell}}.
\]
\excheck $Q_1+Q_2=-q$，
且当 $x\to0$ 时板1感应电荷趋 $-q$；
原题英文正文有一处重复的距离语句，
以其最后明确的“距板1为 $x$”为准。

\electroconductproblem{68}{100}
任意形状导体表面局部面电荷密度为 $\sigma$。
求作用于导体单位表面积上的向外电压力。
\exsolution
在导体外表面，
总电场法向分量
$E_n=\sigma/\varepsilon_0$，
导体内部场为零。
但计算面电荷对自身的净力时，
须排除它的自场；
局部电荷感受到其余电荷产生的平均场
\[
E_{\rm other}=
\frac{E_{\rm outside}+E_{\rm inside}}2
=\frac{\sigma}{2\varepsilon_0}.
\]
所以单位面积受力模
\[
\boxed{P_e=\frac FS=
\sigma E_{\rm other}=
\frac{\sigma^2}{2\varepsilon_0}}.
\]
\excheck $P_e$ 是\emph{正的向外电压力}，
即使 $\sigma<0$，平方后仍为向外张力；
用总场乘 $\sigma$ 会错多出因子2。

\electroconductproblem{69}{101}
一颗半径 $R=1.5\,\mathrm{cm}$ 的带电金属球
总电荷 $q=10\,\mathrm{\mu C}$。
求作用在\emph{其中一个半球表面上的电荷}
所受合电力的模。
\exsolution
球面面电荷密度均匀：
$\sigma=q/(4\pi R^2)$。
第68题给单位面积向外电压力
$P_e=\sigma^2/(2\varepsilon_0)$。
在半球表面，
横向力由对称性抵消，
轴向分量积分
\[
F_z=P_e
\int_0^{2\pi}\!\!\int_0^{\pi/2}
\cos\theta R^2\sin\theta\,\mathrm d\theta\,\mathrm d\phi
=P_e\pi R^2.
\]
因此
\[
\boxed{F=\frac{\sigma^2\pi R^2}{2\varepsilon_0}
=\frac{q^2}{32\pi\varepsilon_0R^2}}.
\]
代入 SI 单位得
\[
\boxed{F\simeq5.0\times10^2\,\mathrm N=0.50\,\mathrm{kN}}.
\]
\excheck 不是简单的 $P_e$ 乘半球表面积 $2\pi R^2$；
因为每个面元受力方向不同，
应投影到共同的对称轴积分。

\electroconductproblem{70}{101}
一个原先中性的导体球，
半径为 $R$，
放在外加均匀电场中后，
其表面出现电荷密度
$\sigma(\theta)=\sigma_0\cos\theta$。
求所有同号诱导电荷所受的合静电力模。
\exsolution
选正半球
$0\le\theta\le\pi/2$，
每处向外电压力为
$P_e(\theta)=\sigma(\theta)^2/(2\varepsilon_0)
=\sigma_0^2\cos^2\theta/(2\varepsilon_0)$。
合力的横向分量抵消，
沿场轴积分：
\[
F=\frac{\sigma_0^2}{2\varepsilon_0}
2\pi R^2\int_0^{\pi/2}
\cos^3\theta\sin\theta\,\mathrm d\theta.
\]
积分为 $1/4$，
故
\[
\boxed{F=\frac{\pi R^2\sigma_0^2}{4\varepsilon_0}}.
\]
\excheck 因只有一个半球同号，
积分限是 $0$ 至 $\pi/2$；
力密度还需投影一次 $\cos\theta$。

\electroconductproblem{71}{101}
强度为
$E=1.0\,\mathrm{kV/cm}$ 的电场
使水分子发生极化。
若此极化等效于每 $N$ 个水分子中
有一个完全沿场方向取向，
求 $N$。
水分子电偶极矩
$p=0.62\times10^{-29}\,\mathrm{C\,m}$。
\exsolution
水的极化强度（线性介质近似）
$P=\varepsilon_0(\varepsilon_r-1)E$。
取室温水的介电常数
$\varepsilon_r\simeq81$、
水的数密度
$n_0=(\rho/M)N_A$，
$\rho\simeq10^3\,\mathrm{kg/m^3}$、
$M\simeq0.018\,\mathrm{kg/mol}$。
若每 $N$ 个分子仅一个完全取向，
等效极化
$P=n_0p/N$。
所以
\[
\boxed{N=\frac{n_0p}
{\varepsilon_0(\varepsilon_r-1)E}}.
\]
代入 $E=10^5\,\mathrm{V/m}$：
\[
\boxed{N\simeq3\times10^3}.
\]
\excheck 原题未明列水的介电常数、密度，
数值须采用室温物性近似；
极化“等效取向分数”不是水分子真正
完全平行排列的比例。

\electroconductproblem{72}{101}
一个不带永久偶极矩、但可极化率为 $\beta$
的非极性分子，距永久偶极矩大小为 $p$
的极性分子 $\ell$，
两个分子连线与 $\bm p$ 同轴。
求两分子间的作用力模。
\exsolution
按本版的\emph{体积极化率}记号，
诱导偶极矩
$\bm p_{\rm ind}=\varepsilon_0\beta\bm E$。
极性分子沿偶极轴的场强为
\[
E(\ell)=\frac{2p}{4\pi\varepsilon_0\ell^3}
=\frac{p}{2\pi\varepsilon_0\ell^3}.
\]
诱导偶极子的势能需包含极化建立功，
故
$U_{\rm ind}=-\frac12\varepsilon_0\beta E^2$。
力的模为
\[
F=\left|-\frac{\mathrm dU_{\rm ind}}{\mathrm d\ell}\right|
=\frac12\varepsilon_0\beta
\left|\frac{\mathrm d(E^2)}{\mathrm d\ell}\right|
=\frac{3\varepsilon_0\beta E^2}{\ell}.
\]
代入轴向电场：
\[
\boxed{F=\frac{3\beta p^2}
{4\pi^2\varepsilon_0\ell^7}}.
\]
\excheck 本题力为吸引。
若使用另一种极化率单位约定
$\bm p_{\rm ind}=\beta_{\rm SI}\bm E$，
则 $\beta_{\rm SI}=\varepsilon_0\beta$，
不能混淆两者。

\electroconductproblem{73}{101}
一个可极化的非极性分子
在半径 $R$、总电荷 $q$ 的均匀带电细圆环轴上，
距环心为 $x$。
求沿轴的作用力为零的位置、
两侧作用力模的极大位置，
并说明 $F_x(x)$ 的图像。
\exsolution
轴上环电场
$E_x=kqx/(R^2+x^2)^{3/2}$。
诱导偶极子的能量
$U=-\frac12\varepsilon_0\beta E_x^2$，
所以
\[
F_x=-\frac{\mathrm dU}{\mathrm dx}
=\frac12\varepsilon_0\beta
\frac{\mathrm d}{\mathrm dx}E_x^2
=\varepsilon_0\beta k^2q^2
\frac{x(R^2-2x^2)}{(R^2+x^2)^4}.
\]
因此当 $x\ge0$ 时，
\[
\boxed{F_x=0\ \Longleftrightarrow\
x=0\ \text{或}\ x=R/\sqrt2}.
\]
在非零位置 $x=R/\sqrt2$，
场强 $E$ 极大而力为零。
设 $u=x/R$，
力的无量纲形状
$f(u)=u(1-2u^2)/(1+u^2)^4$。
求导得到
$f'(u)\propto1-13u^2+10u^4$，
所以两个非零极值位置
\[
\boxed{\frac{x}{R}=
\sqrt{\frac{13-\sqrt{129}}{20}}
\simeq0.29\quad\text{（向外）}},
\]
\[
\boxed{\frac{x}{R}=
\sqrt{\frac{13+\sqrt{129}}{20}}
\simeq1.10\quad\text{（向内）}}.
\]
\excheck $x=0$ 因对称性也是零力点；
原书(a)只列非平凡零力点 $R/\sqrt2$。
力曲线为奇函数，$x>0$ 在
$0<x<R/\sqrt2$ 向外，
更远处向内，并最终按 $x^{-5}$ 衰减：
因圆环远轴场 $E_x\sim kq/x^2$，
$U_{\rm ind}=-\varepsilon_0\beta E_x^2/2\sim-x^{-4}$，
故
$F_x=-\mathrm dU_{\rm ind}/\mathrm dx\sim-x^{-5}$。
这里不能把\emph{点偶极场}的 $x^{-3}$
错当成带电圆环总电荷的远场。

\electroconductproblem{74}{101}
点电荷 $q$ 位于半径为 $R$、
相对介电常数为 $\varepsilon_r$
的均匀各向同性介质球心。
求介质内的极化强度 $\bm P(\bm r)$，
及以球心为中心、半径 $r<R$
的球内\emph{束缚电荷总量}。
\exsolution
由位移矢量高斯定律
$\oint\bm D\cdot\mathrm d\bm S=q$，
\[
\bm D(r)=\frac{q}{4\pi r^2}\bm e_r,\qquad
\bm E=\frac{\bm D}{\varepsilon_0\varepsilon_r}.
\]
极化强度
$\bm P=\varepsilon_0(\varepsilon_r-1)\bm E$，
所以
\[
\boxed{\bm P(\bm r)=
\frac{\varepsilon_r-1}{\varepsilon_r}
\frac{q\,\bm r}{4\pi r^3}}.
\]
束缚电荷净量
$Q'_b=-\oint\bm P\cdot\mathrm d\bm S$，
\[
\boxed{Q'_b=
-q\,\frac{\varepsilon_r-1}{\varepsilon_r}},
\qquad 0<r<R.
\]
\excheck 束缚电荷净量集中在点电荷附近的
分布奇性中；
在 $0<r<R$ 的其余均匀区域，
$\nabla\cdot\bm P=0$。

\electroconductproblem{75}{101}
证明导体与均匀各向同性电介质分界面上，
电介质束缚面电荷密度
$\sigma'_b=-\sigma(\varepsilon_r-1)/\varepsilon_r$，
其中 $\sigma$ 是导体自由电荷面密度。
\exsolution
静电平衡的导体内部 $\bm E=0$，
导体—介质法向位移边界条件给
\[
D_n=\sigma,
\quad D_n=\varepsilon_0\varepsilon_r E_n.
\]
所以介质法向电场
$E_n=\sigma/(\varepsilon_0\varepsilon_r)$。
介质极化为
$\bm P=\varepsilon_0(\varepsilon_r-1)\bm E$，
其在从导体指向介质法向的分量为
$P_n=\sigma(\varepsilon_r-1)/\varepsilon_r$。
但电介质的\emph{外法向}指向导体，
与上述方向相反，
故束缚面电荷
$\sigma'_b=\bm P\cdot\bm n_{\rm diel,out}=-P_n$，
得到
\[
\boxed{\sigma'_b=
-\sigma\frac{\varepsilon_r-1}{\varepsilon_r}}.
\]
\excheck 必须区分“导体外法向”
与“电介质外法向”；
它们在公共界面上方向相反。

\electroconductproblem{76}{101}
一个任意形状导体带净自由电荷 $q$，
其外面包围有一层均匀相对介电常数
$\varepsilon_r$ 的有限厚电介质，
如图3.9。
求介质的内、外两表面上的束缚电荷总量。
\exsolution
导体对周围介质的位移通量为 $q$。
均匀介质中
$\bm P=(\varepsilon_r-1)\bm D/\varepsilon_r$。
内边界按电介质外法向指向导体，
所以
\[
Q'_{\rm inner}=
-\frac{\varepsilon_r-1}{\varepsilon_r}
\oint_{\rm inner}\bm D\cdot\mathrm d\bm S_{\rm conductor}
=\boxed{-q\frac{\varepsilon_r-1}{\varepsilon_r}}.
\]
介质整体无体束缚电荷
（介电常数均匀且介质内部无外加自由体电荷），
$\nabla\cdot\bm P=0$。
所以外表面束缚电荷抵消内表面总量：
\[
\boxed{Q'_{\rm outer}=
+q\frac{\varepsilon_r-1}{\varepsilon_r}}.
\]
\excheck 两表面的总束缚电荷相加为零；
这是有限、均匀且无自由体电荷的介质层。

\electroconductproblem{77}{101}
相对介电常数为 $\varepsilon_r$ 的均匀介质
做成内半径 $a$、外半径 $b>a$ 的球壳。
球壳上存在一定正自由电荷，分别按两种方式：
(a) 均匀分布在内球面；
(b) 均匀分布在球壳体积。
定性给出电场 $E(r)$ 和电势 $\varphi(r)$ 的图线，
并写出分段表达式（总自由电荷记为 $Q$）。
\exsolution
设球壳外部及内腔均为真空，
取无穷远电势为零，
$k=1/(4\pi\varepsilon_0)$。
(a) 自由电荷全在 $r=a$ 球面。
球对称的高斯定律给
\[
\boxed{E_a(r)=
\begin{cases}
0,&r<a,\\
kQ/(\varepsilon_r r^2),&a<r<b,\\
kQ/r^2,&r>b.
\end{cases}}
\]
积分并要求电势连续：
\[
\boxed{\varphi_a(r)=
\begin{cases}
kQ\!\left[\dfrac1{\varepsilon_r a}
+\dfrac{1-1/\varepsilon_r}{b}\right],&r<a,\\
kQ\!\left[\dfrac1{\varepsilon_r r}
+\dfrac{1-1/\varepsilon_r}{b}\right],&a\le r\le b,\\
kQ/r,&r\ge b.
\end{cases}}
\]
(b) 体自由电荷密度
$\rho_f=3Q/[4\pi(b^3-a^3)]$。
球壳内部被包围自由电荷
$Q_{\rm enc}(r)=(4\pi\rho_f/3)(r^3-a^3)$，
故
\[
\boxed{E_b(r)=
\begin{cases}
0,&r<a,\\
\dfrac{\rho_f}{3\varepsilon_0\varepsilon_r}
\left(r-\dfrac{a^3}{r^2}\right),&a\le r\le b,\\
kQ/r^2,&r\ge b.
\end{cases}}
\]
并有
\[
\boxed{\varphi_b(r)=
\begin{cases}
\varphi_b(a),&r<a,\\
\dfrac{kQ}{b}+
\dfrac{\rho_f}{3\varepsilon_0\varepsilon_r}
\left[\dfrac{b^2-r^2}{2}
+a^3\left(\dfrac1b-\dfrac1r\right)\right],
&a\le r\le b,\\
kQ/r,&r\ge b.
\end{cases}}
\]
其中 $\varphi_b(a)$ 取中段 $r=a$ 的值。
图像上：
(a) $E$ 在 $a,b$ 两界面有跳跃，
介质段按 $1/r^2$ 衰减；
(b) $E$ 从 $r=a$ 的零连续上升，
至 $r=b$ 穿出介质时出现介电跳跃，
再按 $1/r^2$ 衰减；
两种情况下 $\varphi(r)$ 均连续，
且内腔为常势。
\excheck 介质边界若无自由面电荷，
连续的是 $D_n$，
不是 $E_n$；而在(a)的内球面
存在自由面电荷，$D_n$ 必有相应跃变。
% === END 17-electro-conductors-066-077.tex ===

% === BEGIN 17-electro-conductors-078-088.tex blob 44e0f4b65b6d21696d264a9c8a0775adad2d7ca7 ===
% Irodov Part III §3.2, original 78--88, PDF pp.111--112,
% printed pp.102--103; Fig 3.10--3.14 checked against source images.
\electroconductproblem{78}{102}
玻璃与真空分界面上的点 A
（见图3.10），真空侧电场强度
$E_0=10.0\,\mathrm{V/m}$，
与从玻璃指向真空的法线夹角
$\alpha_0=30^\circ$。
求玻璃内近边界处的场强 $E$、
与法线夹角 $\alpha$、
以及玻璃的束缚表面电荷密度。
\exsolution
设玻璃相对介电常数为 $\varepsilon_r$。
界面无自由面电荷时，
法向 $\bm D$ 连续、切向 $\bm E$ 连续：
\[
E_n=\frac{E_0\cos\alpha_0}{\varepsilon_r},
\qquad E_t=E_0\sin\alpha_0.
\]
因此
\[
\boxed{E=E_0\sqrt{
\frac{\cos^2\alpha_0}{\varepsilon_r^2}
+\sin^2\alpha_0}},
\qquad
\boxed{\tan\alpha=\varepsilon_r\tan\alpha_0}.
\]
玻璃极化
$\bm P=\varepsilon_0(\varepsilon_r-1)\bm E$，
故面束缚电荷
\[
\boxed{\sigma_b=
\varepsilon_0\frac{\varepsilon_r-1}{\varepsilon_r}
E_0\cos\alpha_0}.
\]
取原书数值答案所用的玻璃介电常数
$\varepsilon_r\simeq6$，
得到
\[
\boxed{E\simeq5.2\,\mathrm{V/m}},\quad
\boxed{\alpha\simeq74^\circ},\quad
\boxed{\sigma_b\simeq64\,\mathrm{pC/m^2}}.
\]
\excheck 玻璃的 $\varepsilon_r$ 未在本题直接列数值，
因此数值须说明采用的常用玻璃介电常数，
公式不依赖这一选择。

\electroconductproblem{79}{102}
一均匀各向同性电介质半空间的表面为平面，
其相对介电常数为 $\varepsilon_r$。
真空侧近表面的均匀电场
$\bm E_0$ 与法线成 $\theta$ 角
（图3.11）。
假设两侧场分别均匀，求：
(a) 以界面上点为球心、半径 $R$ 的球面的
$\bm E$ 通量；
(b) 沿图中与 $\bm E_0$ 平行的法截面内、
跨越界面的长为 $\ell$ 的狭矩形闭路的
$\bm D$ 环流。
\exsolution
无自由面电荷时
$E_{{\rm diel},n}=E_0\cos\theta/\varepsilon_r$，
$E_{{\rm diel},t}=E_0\sin\theta$。
(a) 球面的一半处于真空、一半在电介质。
每半球均匀场的通量由投影面积
$\pi R^2$ 给出，
所以
\[
\Phi_E=
\pi R^2E_0\cos\theta-
\pi R^2\frac{E_0\cos\theta}{\varepsilon_r}.
\]
即
\[
\boxed{\Phi_E=
\frac{\varepsilon_r-1}{\varepsilon_r}
\pi R^2 E_0\cos\theta}.
\]
(b) 让狭矩形两个长边分别位于界面的
真空侧和介质侧，
并按图3.11给定方向取闭路。
切向电场连续，
但切向位移分别为
$D_{0,t}=\varepsilon_0E_0\sin\theta$、
$D_{{\rm diel},t}=\varepsilon_0\varepsilon_rE_0\sin\theta$。
跨界面短边在宽度趋零时不贡献，
所以
\[
\boxed{\oint_\Gamma\bm D\cdot\mathrm d\bm l
=-\varepsilon_0(\varepsilon_r-1)
\ell E_0\sin\theta}.
\]
\excheck $\bm E$ 满足无旋性，
但 $\bm D$ 穿越不同介电常数的边界时
其切向分量一般不连续，
故 $\oint\bm D\cdot\mathrm d\bm l$
可以非零。

\electroconductproblem{80}{102}
一块沿 $x$ 轴法向、厚度为 $2d$
的无限大均匀电介质平板，
相对介电常数为 $\varepsilon_r$，
在 $|x|<d$ 内存在均匀\emph{自由体电荷密度}
$\rho_f$。
设中心电势为零，求：
(a) 电场 $E_x(x)$、电势 $\varphi(x)$；
(b) 束缚电荷的体密度和表面密度。
\exsolution
系统关于 $x=0$ 对称，
$\bm D$ 在介质内部满足
$\mathrm dD_x/\mathrm dx=\rho_f$，
由 $D_x(0)=0$，
$D_x=\rho_fx$。
在外侧
$D_x=\rho_fd\,\operatorname{sgn}x$。
因此
\[
\boxed{E_x(x)=
\begin{cases}
\dfrac{\rho_f x}{\varepsilon_0\varepsilon_r},
&|x|\le d,\\[2mm]
\dfrac{\rho_fd}{\varepsilon_0}\operatorname{sgn}x,
&|x|\ge d.
\end{cases}}
\]
用 $E_x=-\varphi'(x)$，
并取 $\varphi(0)=0$、电势连续，
\[
\boxed{\varphi(x)=
\begin{cases}
-\dfrac{\rho_fx^2}{2\varepsilon_0\varepsilon_r},
&|x|\le d,\\[2mm]
-\dfrac{\rho_fd^2}{2\varepsilon_0\varepsilon_r}
-\dfrac{\rho_fd}{\varepsilon_0}(|x|-d),
&|x|\ge d.
\end{cases}}
\]
极化
$\bm P=(\varepsilon_r-1)\bm D/\varepsilon_r$。
内部束缚体密度
\[
\boxed{\rho'_b=-\nabla\cdot\bm P
=-\rho_f\frac{\varepsilon_r-1}{\varepsilon_r}}.
\]
两侧束缚面电荷均同号且大小相等：
\[
\boxed{\sigma'_b(\pm d)=
+\rho_fd\frac{\varepsilon_r-1}{\varepsilon_r}}.
\]
\excheck 电场为\emph{奇函数}，
电势为\emph{偶函数}；
两侧束缚面电荷与内部束缚体电荷
总量每单位面积相消：
$2\sigma'_b+2d\rho'_b=0$。

\electroconductproblem{81}{102}
半径 $R$ 的均匀各向同性电介质实心球
相对介电常数为 $\varepsilon_r$，
内部自由体电荷密度
$\rho_f>0$ 均匀分布。
求 (a) 内外电场与电势的半径函数；
(b) 束缚电荷体密度与表面密度；
说明图线形状。
\exsolution
高斯定律对位移矢量：
$D_r(r)=\rho_fr/3$（$r<R$），
$D_r(r)=\rho_fR^3/(3r^2)$（$r>R$）。
因此
\[
\boxed{E(r)=
\begin{cases}
\dfrac{\rho_fr}{3\varepsilon_0\varepsilon_r},
&r<R,\\[2mm]
\dfrac{\rho_fR^3}{3\varepsilon_0r^2},
&r>R.
\end{cases}}
\]
取无穷远电势为零，
球面
$\varphi(R)=\rho_fR^2/(3\varepsilon_0)$，
由内场积分得到
\[
\boxed{\varphi(r)=
\begin{cases}
\dfrac{\rho_fR^2}{3\varepsilon_0}
+\dfrac{\rho_f(R^2-r^2)}
{6\varepsilon_0\varepsilon_r},&r\le R,\\[2mm]
\dfrac{\rho_fR^3}{3\varepsilon_0r},&r\ge R.
\end{cases}}
\]
内部极化
$\bm P=(\varepsilon_r-1)\rho_f\bm r/(3\varepsilon_r)$，
所以
\[
\boxed{\rho'_b=
-\rho_f\frac{\varepsilon_r-1}{\varepsilon_r}},
\quad
\boxed{\sigma'_b=
\frac{\varepsilon_r-1}{\varepsilon_r}
\frac{\rho_fR}{3}}.
\]
\excheck $E(r)$ 在球内线性增加，
在 $r=R$ 因介电常数改变而跳跃，
球外按 $1/r^2$ 衰减；
电势连续，内段为抛物线、外段为 $1/r$。
全部束缚体、面电荷代数和为零。

\electroconductproblem{82}{102}
半径 $R$、厚度 $d\ll R$ 的薄圆盘形电介质，
具有均匀固定极化强度 $\bm P$，
其方向\emph{位于圆盘平面内}。
求圆盘中心电场。
\exsolution
取 $\bm P=P\bm e_x$。
上下两个平面因其法线与 $\bm P$ 垂直，
束缚面电荷为零；
只有侧面 $r=R$ 带束缚电荷
$\sigma_b=P\cos\phi$。
在圆盘中心，侧面元
$\mathrm dS=R\,\mathrm d\phi\,\mathrm dz$
产生
$\mathrm d\bm E=
-k\sigma_b(R\bm e_r+z\bm e_z)
\,\mathrm dS/(R^2+z^2)^{3/2}$。
积分侧壁
$z\in[-d/2,d/2]$ 及角度，
由
$\int_0^{2\pi}\cos^2\phi\,\mathrm d\phi=\pi$，
得
\[
\bm E(0)=
-\frac{PR^2}{4\varepsilon_0}
\left(\int_{-d/2}^{d/2}
\frac{\mathrm dz}{(R^2+z^2)^{3/2}}\right)\bm e_x.
\]
在 $d\ll R$ 时，
括号内约为 $d/R^3$，
故
\[
\boxed{\bm E(0)\simeq
-\frac{d}{4\varepsilon_0R}\bm P}.
\]
\excheck 圆盘极化平行盘面，
故只有侧壁带束缚电荷；
若误当成盘面法向极化，将得到完全不同的场。

\electroconductproblem{83}{102--103}
无自由电荷的无限大电介质平板，
厚度为 $2d$，取中面为 $x=0$。
其极化分布
$\bm P(x)=\bm P_0(1-x^2/d^2)$，
$\bm P_0$ 垂直平板。
求板内电场与两表面间电势差。
\exsolution
设 $\bm P_0=P_0\bm e_x$。
没有自由电荷，
$\mathrm dD_x/\mathrm dx=0$。
且极化在 $|x|=d$ 处为零，
无外加电场的孤立板在外侧场为零，
故 $D_x=0$，即
$\varepsilon_0E_x+P_x=0$。
因此
\[
\boxed{\bm E(x)=
-\frac{\bm P_0}{\varepsilon_0}
\left(1-\frac{x^2}{d^2}\right)},\qquad |x|\le d.
\]
两表面间电势差定义为
$\varphi(d)-\varphi(-d)$，
有
\[
\Delta\varphi=
-\int_{-d}^dE_x\,\mathrm dx
=\frac{P_0}{\varepsilon_0}
\int_{-d}^d(1-x^2/d^2)\,\mathrm dx,
\]
故
\[
\boxed{\Delta\varphi=\frac{4dP_0}{3\varepsilon_0}}.
\]
\excheck 电介质平板外侧场为零是此\emph{无外加场}
的边界条件；否则 $D$ 可能带另一常数。

\electroconductproblem{84}{103}
两板间初始为空气的平行板电容器，
间距 $d$，初电场 $E_0$。
现按图3.12把\emph{沿场强方向的半层厚度}
替换为相对介电常数 $\varepsilon_r$ 的材料，
形成两个各厚 $d/2$ 的\emph{串联层}。
求气隙部分1和介质部分2的
$\bm E,\bm D$ 模：
(a) 极板电压不变；(b) 极板自由电荷不变。
\exsolution
界面没有自由面电荷，故
$D_1=D_2=D$。
两层法向场
$E_1=D/\varepsilon_0$、
$E_2=D/(\varepsilon_0\varepsilon_r)$。
(a) 电压固定为初值 $V_0=E_0d$，
\[
V_0=\frac d2(E_1+E_2)
=\frac{dD}{2\varepsilon_0}
(1+1/\varepsilon_r).
\]
所以
\[
\boxed{E_1=\frac{2\varepsilon_r}{\varepsilon_r+1}E_0},
\qquad
\boxed{E_2=\frac2{\varepsilon_r+1}E_0},
\quad
\boxed{D_1=D_2=
\frac{2\varepsilon_r}{\varepsilon_r+1}\varepsilon_0E_0}.
\]
(b) 极板自由电荷不变，
$D=D_0=\varepsilon_0E_0$，因此
\[
\boxed{E_1=E_0,\quad E_2=E_0/\varepsilon_r},
\qquad
\boxed{D_1=D_2=\varepsilon_0E_0}.
\]
\excheck 图3.12 是\emph{沿电场方向的串联半层}，
与第85题的并联半面积构型不同。

\electroconductproblem{85}{103}
把前题的电介质改按图3.13布置，
占电容器极板面积的一半，
且每块区域均跨越\emph{全板间距}，
形成空气区1与介质区2的\emph{并联}。
其余条件相同。求 (a) 电压不变；
(b) 极板电荷不变时各区 $E,D$。
\exsolution
忽略边缘效应时，
两区域上下极板相同，
因此具有相同电压 $V$，
$E_1=E_2=V/d$。
(a) $V=V_0=E_0d$，
故
\[
\boxed{E_1=E_2=E_0},\qquad
\boxed{D_1=\varepsilon_0E_0,\quad
D_2=\varepsilon_r\varepsilon_0E_0}.
\]
(b) 极板总面积为 $S$，
原自由电荷
$Q_0=\varepsilon_0E_0S$。
新电容为两个面积 $S/2$ 的电容并联：
\[
C_{\rm new}=\frac{\varepsilon_0S}{2d}
(1+\varepsilon_r).
\]
因此
$V=Q_0/C_{\rm new}
=2E_0d/(1+\varepsilon_r)$，
\[
\boxed{E_1=E_2=\frac{2E_0}{1+\varepsilon_r}},
\]
\[
\boxed{D_1=\frac{2\varepsilon_0E_0}{1+\varepsilon_r},
\quad
D_2=\frac{2\varepsilon_r\varepsilon_0E_0}
{1+\varepsilon_r}}.
\]
\excheck 固定\emph{总电荷}不表示两半极板
各自的自由电荷密度保持原值；
并联两区的自由电荷会重新分配。

\electroconductproblem{86}{103}
一个同心球形电容器，
内外导体球壳同心、内导体带自由电荷 $q$。
两极板之间的空间按图3.14
沿穿过球心的平面划为两半：
一半是真空，一半填相对介电常数
$\varepsilon_r$ 的均匀材料。
求极板之间的场强模 $E(r)$。
\exsolution
球形等势导体的边界迫使电场沿径向。
两半空间共用同一径向电势梯度，
故两侧 $E_r(r)$ 相同，
而 $D$ 不同：
$D_{\rm vac}=\varepsilon_0 E_r$，
$D_{\rm diel}=\varepsilon_0\varepsilon_r E_r$。
对半径 $r$ 的球面取高斯通量，
真空和介质各占半球面积 $2\pi r^2$：
\[
q=\oint\bm D\cdot\mathrm d\bm S
=2\pi r^2\varepsilon_0E_r+
2\pi r^2\varepsilon_r\varepsilon_0E_r.
\]
因此
\[
\boxed{E(r)=
\frac{|q|}{2\pi\varepsilon_0(1+\varepsilon_r)r^2}}.
\]
\excheck 两半分区\emph{并联}，不能把介质
误当成径向串联层；
当 $\varepsilon_r=1$ 恢复
$|q|/(4\pi\varepsilon_0r^2)$。

\electroconductproblem{87}{103}
两个带同号等量电荷的相同小球，
由等长绝缘线悬挂于同一点。
将周围介质由空气换成煤油后，
两线间张角\emph{保持不变}。
求小球材料密度 $\rho$，
煤油密度记为 $\rho_0$，
相对介电常数为 $\varepsilon_r$。
\exsolution
在空气中静电排斥力为 $F_e$，
球的有效重力约为 $mg$。
在煤油中，同样球间距离下
库仑力降为 $F_e/\varepsilon_r$；
与此同时浮力使有效重力变为
\[
mg-\rho_0Vg=(\rho-\rho_0)Vg
=mg\left(1-\frac{\rho_0}{\rho}\right).
\]
张角不变要求
\[
\frac{F_e}{mg}
=\frac{F_e/\varepsilon_r}
{mg(1-\rho_0/\rho)}.
\]
约去公共因子，
$1-\rho_0/\rho=1/\varepsilon_r$，
得
\[
\boxed{\rho=
\frac{\varepsilon_r}{\varepsilon_r-1}\rho_0}.
\]
取煤油常用近似
$\varepsilon_r\simeq2.0$、
$\rho_0\simeq0.80\,\mathrm{g/cm^3}$，
\[
\boxed{\rho\simeq1.6\,\mathrm{g/cm^3}}.
\]
\excheck 不能只考虑介电常数使库仑力减小，
必须同时计入煤油对小球的浮力。

\electroconductproblem{88}{103}
一个半径
$R=3.0\,\mathrm{cm}$ 的均匀各向同性电介质球，
相对介电常数为
$\varepsilon_r=5.00$，
球内部产生均匀电场
$E=100\,\mathrm{V/m}$。
求：(a) 最大束缚面电荷密度；
(b) 同号束缚电荷的总量。
\exsolution
线性介质的极化矢量为
$\bm P=\varepsilon_0(\varepsilon_r-1)\bm E$。
均匀极化球的表面束缚电荷
$\sigma_b(\theta)=\bm P\cdot\bm n
=P\cos\theta$。
故最大面密度位于极化轴两极：
\[
\boxed{\sigma_{b,\max}=
\varepsilon_0(\varepsilon_r-1)E}.
\]
代入数值：
\[
\boxed{\sigma_{b,\max}\simeq
(8.854\times10^{-12})(4)(100)
=3.5\,\mathrm{nC/m^2}}.
\]
同号电荷全部位于一半球面，
其总量是
\[
Q'_+=\int_0^{2\pi}\int_0^{\pi/2}
P\cos\theta R^2\sin\theta
\,\mathrm d\theta\,\mathrm d\phi
=\pi R^2P.
\]
因此
\[
\boxed{Q'_+=\pi R^2
\varepsilon_0(\varepsilon_r-1)E
\simeq1.0\times10^{-11}\,\mathrm C
=10\,\mathrm{pC}}.
\]
\excheck 负束缚电荷总量为 $-Q'_+$，
所以中性电介质球的束缚电荷代数和为零。
% === END 17-electro-conductors-078-088.tex ===

% === BEGIN 17-electro-conductors-089-100.tex blob bd24c63a844884259673ad7f5acb786c6d669aaa ===
% Irodov Part III §3.2 (89--100), user PDF pp.112--114,
% printed pp.103--105. Independent Part-III numbering.
% Conventions: eps=relative permittivity; r_perp=distance from interface normal.
\electroconductproblem{89}{103}
真空半空间内有点电荷 $q$，距平面电介质界面为
$\ell$。介质占据另一半无限空间，相对介电常数为
$\varepsilon$。求：(a) 距电荷在界面上投影点
的径向距离为 $s$ 处，界面束缚电荷面密度；
讨论 $\ell\to0$；(b) 界面总束缚电荷。
\exsolution
取界面 $z=0$，真电荷位于 $(0,0,\ell)$。
介质边界的镜像法给真空区等效镜像电荷
$q'=q(1-\varepsilon)/(1+\varepsilon)$，
位于 $z=-\ell$。介质区的等效场相当于
$2q/(1+\varepsilon)$ 位于真电荷原位置。
介质侧的法向极化强度为
$\bm P=\varepsilon_0(\varepsilon-1)\bm E_{\rm diel}$。
在介质外法向为 $+\bm e_z$ 的边界，
介质内 $E_z=-[2q/(1+\varepsilon)]
\ell/[4\pi\varepsilon_0(s^2+\ell^2)^{3/2}]$，
于是
\[
\boxed{\sigma_b(s)=
-\frac{q\ell(\varepsilon-1)}
{2\pi(\varepsilon+1)(s^2+\ell^2)^{3/2}}}.
\]
固定 $s>0$ 时 $\ell\to0$ 有 $\sigma_b\to0$；
但在 $s=0$ 处呈集中奇性，
所以不能将总电荷误断为零。
积分平面：
\[
Q_b=\int_0^\infty2\pi s\,\sigma_b(s)\,\mathrm ds
=-\frac{q(\varepsilon-1)}{\varepsilon+1}
\ell\int_0^\infty\frac{s\,\mathrm ds}{(s^2+\ell^2)^{3/2}}
=\boxed{-q\frac{\varepsilon-1}{\varepsilon+1}}.
\]
\excheck 作为 $\varepsilon\to\infty$ 的导体极限，
$Q_b\to-q$，符合接地导体镜像解。

\electroconductproblem{90}{103}
利用上一题结果，求电介质界面束缚电荷对
位于界面上方 $\ell$ 处点电荷 $q$ 的作用力。
\exsolution
真空一侧的介质界面作用，可由位于镜像位置的
$q'=q(1-\varepsilon)/(1+\varepsilon)$
来表示，镜像距真电荷 $2\ell$。
因此作用力沿界面法线：
\[
F_z=\frac1{4\pi\varepsilon_0}
\frac{q q'}{(2\ell)^2}
=-\frac{q^2(\varepsilon-1)}
{16\pi\varepsilon_0\ell^2(\varepsilon+1)}.
\]
其模
\[
\boxed{F=\frac{q^2(\varepsilon-1)}
{16\pi\varepsilon_0\ell^2(\varepsilon+1)}}.
\]
当 $\varepsilon>1$，力指向介质。
\excheck $\varepsilon=1$ 时介质与真空不可区分，
力为零；$\varepsilon\to\infty$ 恢复导体镜像力。

\electroconductproblem{91}{104}
点电荷 $q$ 位于真空和均匀无限电介质
（相对介电常数 $\varepsilon$）的平面分界线上。
求两侧距电荷 $r$ 处的
$\bm D,\bm E$ 模及电势 $\varphi(r)$。
\exsolution
由界面法向位移与切向电场边界条件，
平面两侧的径向电场可取
$E(r)=A/r^2$，两侧取同一系数。
以电荷为球心的高斯球在真空、介质
各有一半球面积 $2\pi r^2$：
\[
q=2\pi r^2(\varepsilon_0E+
\varepsilon_0\varepsilon E)
=2\pi\varepsilon_0(1+\varepsilon)Ar^0.
\]
因此
\[
\boxed{E_{\rm vac}=E_{\rm diel}
=\frac{|q|}{2\pi\varepsilon_0(1+\varepsilon)r^2}}.
\]
位移矢量模为
\[
\boxed{D_{\rm vac}=\frac{|q|}
{2\pi(1+\varepsilon)r^2}},\qquad
\boxed{D_{\rm diel}=\frac{\varepsilon|q|}
{2\pi(1+\varepsilon)r^2}}.
\]
取无穷远电势零，径向积分得
\[
\boxed{\varphi(r)=
\frac{q}{2\pi\varepsilon_0(1+\varepsilon)r}}.
\]
\excheck 电势与 $\bm E$ 在两侧连续，
而 $\bm D$ 的模因介电常数不同而不同。

\electroconductproblem{92}{104}
在均匀电介质半空间内，距与真空分界面
$\ell$ 处有一个小导体球，带电 $q$。
介质相对介电常数 $\varepsilon$，
球半径远小于 $\ell$。
求平面束缚电荷密度的分布及 $\ell\to0$ 的行为。
\exsolution
在点电荷近似下，真电荷位于介质侧
$z=-\ell$。介质侧镜像电荷
$q'=q(\varepsilon-1)/(\varepsilon+1)$
位于 $z=+\ell$。
介质中的场为
$\bm E=(4\pi\varepsilon_0\varepsilon)^{-1}$
乘两点电荷叠加。
在界面 $z=0^-$，
两电荷的法向场为
\[
E_z=\frac{(q-q')\ell}
{4\pi\varepsilon_0\varepsilon(s^2+\ell^2)^{3/2}}
=\frac{q\ell}
{2\pi\varepsilon_0\varepsilon(\varepsilon+1)
(s^2+\ell^2)^{3/2}}.
\]
介质表面外法向为 $+\bm e_z$，所以
\[
\boxed{\sigma_b(s)=
\varepsilon_0(\varepsilon-1)E_z
=\frac{q\ell(\varepsilon-1)}
{2\pi\varepsilon(\varepsilon+1)
(s^2+\ell^2)^{3/2}}}.
\]
\excheck 对固定 $s>0$，$\ell\to0$ 时密度趋零；
当 $\ell$ 过小接近球半径时，点电荷近似本身失效。

\electroconductproblem{93}{104}
均匀电介质占据导体平面一侧半空间，
相对介电常数 $\varepsilon$。
距接地导体平面 $\ell$ 处有小金属球，
带电 $q$，球半径可忽略。
求导体边界处的介质束缚面电荷密度。
\exsolution
将导体平面视为 $z=0$，
真电荷位于 $z=\ell$，
镜像 $-q$ 位于 $z=-\ell$。
介质内部的法向场
$E_z(0^+)=-q\ell/
[2\pi\varepsilon_0\varepsilon(s^2+\ell^2)^{3/2}]$。
介质的外法向在底部是 $-\bm e_z$，
故
\[
\sigma_b=\bm P\cdot(-\bm e_z)
=-\varepsilon_0(\varepsilon-1)E_z(0^+).
\]
于是
\[
\boxed{\sigma_b(s)=
\frac{q\ell(\varepsilon-1)}
{2\pi\varepsilon(s^2+\ell^2)^{3/2}}}.
\]
\excheck 对正 $q$，靠近导体的介质束缚面
带正电，但导体感应\emph{自由电荷}为负电；
二者不能混为一谈。

\electroconductproblem{94}{104}
图3.15：两平行导体板短接，板距 $d$，
板间放一块沿板法向均匀固定极化的介质板，
厚度 $h<d$，极化强度为 $\bm P$。
求介质内部与剩余气隙中的 $\bm E,\bm D$。
\exsolution
取极化方向为正法向，用标量 $P$ 表示。
两个区域都没有自由体电荷，
法向位移恒为 $D$。
在介质板内部
$D=\varepsilon_0E_{\rm diel}+P$，
气隙内
$D=\varepsilon_0E_{\rm gap}$。
由于两极板短接，总电压为零：
\[
(d-h)E_{\rm gap}+hE_{\rm diel}=0.
\]
代入 $E_{\rm gap}=D/\varepsilon_0$、
$E_{\rm diel}=(D-P)/\varepsilon_0$ 得
$dD-hP=0$，因此
\[
\boxed{\bm D=\frac hd\bm P
\quad\text{（两区域）}},
\]
\[
\boxed{\bm E_{\rm gap}=\frac{h}{\varepsilon_0d}\bm P},
\qquad
\boxed{\bm E_{\rm diel}=
-\frac{1-h/d}{\varepsilon_0}\bm P}.
\]
\excheck $h=d$ 时介质完全填满，
气隙消失、内部电场为零；
$h\to0$ 时电介质作用趋零。

\electroconductproblem{95}{104}
无限长圆柱体电介质沿径向具有固定极化
$\bm P=\alpha\bm r_\perp$，
$\alpha>0$。求内部束缚体电荷密度。
\exsolution
束缚体电荷密度定义为
$\rho_b=-\nabla\cdot\bm P$。
在柱坐标里，
$P_r=\alpha r$，
\[
\nabla\cdot\bm P=
\frac1r\frac{\partial(rP_r)}{\partial r}
=\frac1r\frac{\partial(\alpha r^2)}{\partial r}
=2\alpha.
\]
所以
\[
\boxed{\rho_b=-2\alpha}.
\]
\excheck $\alpha$ 的单位为
$\mathrm{C/m^3}$，密度与半径无关；
若圆柱半径为 $R$，
侧表面还带束缚面电荷
$\sigma_b=\alpha R$。

\electroconductproblem{96}{104}
均匀固定极化的介质球，极化为 $\bm P$。
(a) 利用正负电荷的微小错位证明球内电场；
(b) 证明球外是中心点偶极子的电场，
并求其偶极矩。
\exsolution
极化均匀，束缚体密度
$\rho_b=-\nabla\cdot\bm P=0$，
球表面密度 $\sigma_b=P\cos\theta$。
其可看作密度为 $+\rho,-\rho$ 的两均匀带电球
相对微小偏移 $\bm\delta$，满足
$\rho\bm\delta=\bm P$。
在两球共同内部，场为
\[
\bm E=
\frac{\rho}{3\varepsilon_0}\bm r
-\frac{\rho}{3\varepsilon_0}(\bm r+\bm\delta)
=\boxed{-\frac{\bm P}{3\varepsilon_0}}.
\]
球外总电荷零，最低非零多极矩为偶极矩
\[
\boxed{\bm p=V\bm P=
\frac{4\pi R^3}{3}\bm P}.
\]
其球外电势可直接由
$\sigma_b=P\cos\theta$ 的球面积分验证为
\[
\boxed{\varphi(\bm r)=
\frac{\bm p\cdot\bm r}
{4\pi\varepsilon_0r^3}},\qquad r>R.
\]
\excheck 球内均匀退极化场指向
$-\bm P$；球外场不等于同一均匀场。

\electroconductproblem{97}{104}
无限均匀固定极化的电介质内部挖去
一个球形空腔。介质极化强度为 $\bm P$，
远离空腔的实际电场为 $\bm E_\infty$。
求空腔内电场 $\bm E_{\rm cav}$。
\exsolution
线性叠加：用一个均匀极化为 $-\bm P$ 的球
叠加到原均匀极化介质，
刚好抵消球内的极化，即挖去空腔。
根据第96题，一颗均匀极化为
$\bm P_{\rm sph}$ 的球在内部产生
$-\bm P_{\rm sph}/(3\varepsilon_0)$。
因此空腔附加场
$\Delta\bm E=-(-\bm P)/(3\varepsilon_0)
=+\bm P/(3\varepsilon_0)$，
所以
\[
\boxed{\bm E_{\rm cav}=
\bm E_\infty+\frac{\bm P}{3\varepsilon_0}}.
\]
\excheck 注意本题的远场 $\bm E_\infty$
是\emph{原介质中的实际宏观场}。
若原书简答出现相反符号，可能采用不同的
场变量约定；本式以叠加法的极化定义为准。

\electroconductproblem{98}{104}
均匀相对介电常数为 $\varepsilon$ 的
电介质球置于外加均匀场 $\bm E_0$，
球在作用下均匀极化。
利用第96题求球内电场与极化强度。
\exsolution
球内总电场为施加场与自身退极化场之和：
\[
\bm E=\bm E_0-\frac{\bm P}{3\varepsilon_0}.
\]
线性介质满足
$\bm P=\varepsilon_0(\varepsilon-1)\bm E$，
代入：
\[
\bm E\left(1+\frac{\varepsilon-1}{3}\right)
=\bm E_0.
\]
故
\[
\boxed{\bm E=\frac3{\varepsilon+2}\bm E_0},
\quad
\boxed{\bm P=
\frac{3\varepsilon_0(\varepsilon-1)}
{\varepsilon+2}\bm E_0}.
\]
\excheck $\varepsilon=1$ 时极化为零，
球内场等于外加场；
$\varepsilon\to\infty$ 时
球内场趋零。

\electroconductproblem{99}{105}
无限长圆柱形介质固定均匀极化，
极化 $\bm P$ 垂直于圆柱轴线。
求圆柱内部电场。
\exsolution
取圆柱轴沿 $z$、极化沿 $+x$，
侧表面束缚电荷密度
$\sigma_b=P\cos\phi$，内部无束缚体电荷。
二维边值问题的内外电势可以取
\[
\varphi_{\rm in}=Ar\cos\phi,\qquad
\varphi_{\rm out}=B\frac{\cos\phi}{r},
\]
其中表面电势连续给 $B=AR^2$。
在 $r=R$，法向场跃变
$E_{{\rm out},n}-E_{{\rm in},n}=
\sigma_b/\varepsilon_0$。
由
$E_{{\rm in},n}=-A\cos\phi$、
$E_{{\rm out},n}=+A\cos\phi$，
得 $2A=P/\varepsilon_0$。
因此
\[
\boxed{\bm E_{\rm in}=
-\frac{\bm P}{2\varepsilon_0}}.
\]
\excheck 与球体 $1/3$ 系数不同，
圆柱横向退极化因子是 $1/2$。

\electroconductproblem{100}{105}
无限长均匀介质圆柱置于均匀外场
$\bm E_0$ 内，圆柱轴线与 $\bm E_0$ 垂直。
圆柱相对介电常数为 $\varepsilon$。
利用上一题结果求柱内场与极化强度。
\exsolution
根据第99题，均匀极化产生
$-\bm P/(2\varepsilon_0)$ 退极化场，
因此
\[
\bm E=\bm E_0-\frac{\bm P}{2\varepsilon_0}.
\]
材料本构关系
$\bm P=\varepsilon_0(\varepsilon-1)\bm E$，
代入：
\[
\left(1+\frac{\varepsilon-1}{2}\right)\bm E
=\bm E_0,
\]
从而
\[
\boxed{\bm E=
\frac2{\varepsilon+1}\bm E_0},
\qquad
\boxed{\bm P=
\frac{2\varepsilon_0(\varepsilon-1)}
{\varepsilon+1}\bm E_0}.
\]
\excheck 当介质常数趋于1时，
内部场还原为外加场；
$\varepsilon\to\infty$ 时内部电场趋零。
% === END 17-electro-conductors-089-100.tex ===
